% AI-produced by GPT-6 Astra, OpenAI, 20 November 2026.
% Candidate exact normalization of the A5.019 finite-pool gate.
% Not included in the compiled paper; no closure or movement claimed.
\subsection*{Candidate: a standard angular Lipschitz bound in the exact A5 gate}

\textbf{Proposition.}\quad
Fix the named-coefficient complete theory and the finitely many admissible
radial models of A5.019. On one such model let $W$ be the scalar-certificate
parameter set, let $L$ be its fixed nonempty compact polyhedron, let
$\Sigma=\partial[-1,1]^m$ with $m\ge1$, and let
\[
 \Theta_w:(0,1]\times L\times\Sigma\longrightarrow\Sigma
 \tag{AL.1}
\]
be the definable angular map. For each fixed $w$ it is jointly continuous
in $(t,z,u)$; no joint continuity in $w$ is assumed.
There is one fixed definable family
\[
 r_{w,t,z}:\Sigma\longrightarrow\Sigma
 \tag{AL.2}
\]
of homeomorphisms, jointly continuous with its inverse in $(t,z,u)$ for
every fixed $w$, and an ordinary integer $K\ge1$, such that
\[
 \begin{split}
 \|r_{w,t,z}(u)-r_{w,t,z}(v)\|_\infty&\le K\|u-v\|_\infty,\\
 \|\widehat\Theta_w(t,z,u)-\widehat\Theta_w(t,z,v)\|_\infty
                                      &\le K\|u-v\|_\infty,
 \end{split}
 \tag{AL.3}
\]
where
\[
 \widehat\Theta_w(t,z,u)=\Theta_w(t,z,r_{w,t,z}(u)).
 \tag{AL.4}
\]
The same fixed formulas and the same ordinary integer bound work in
every model of the fixed complete theory. A single $K$ may be used for
all the finitely many original radial models.

Put
\[
 \chi_{w,t}(z,u)=(z,r_{w,t,z}(u)).
 \tag{AL.5}
\]
This is a homeomorphism of $L\times\Sigma$ preserving
$\operatorname{pr}_L$ exactly. The straightener correspondence
\[
 \widehat A_t=\chi_{w,t}^{-1}A_t,
 \qquad A_t=\chi_{w,t}\widehat A_t
 \tag{AL.6}
\]
preserves the exact angular equation and gives the exact error identity
\[
 \sup_{z,u}\|\operatorname{pr}_L\widehat A_t(z,u)-z\|
 =\sup_{z,u}\|\operatorname{pr}_L A_t(z,u)-z\|.
 \tag{AL.7}
\]
Consequently A5.019's all-model finite-pool assertion is equivalent to
the same assertion for this finite collection of uniformly angular
$K$-Lipschitz maps $\widehat\Theta$.
The bound in (AL.3) concerns the known angular input and the known gauge;
it does not bound the unknown straighteners or their formula complexity.

\begin{proof}
\textbf{A fixed cubical model and a bounded extension.}\quad
Use a rational finite triangulation of $\Sigma$. The explicit cubical
PL model in A154 sends the barycentre of each face to its characteristic
vector and extends affinely on chains of faces. It gives a fixed PL
homeomorphism
\[
 b:\Sigma\longrightarrow Q\subset C=[0,1]^v,
\]
where $Q$ is a finite union of coordinate faces of $C$. Both $b$ and its
inverse have Lipschitz bounds by ordinary integers: they have fixed
finite rational affine descriptors. The case $m=1$ is included; then
$\Sigma$ is a finite zero-dimensional polyhedron.

For the formal parameter $a=(w,t,z)$, define
\[
 F_a(x)=\tfrac12\bigl(\Theta_w(t,z,b^{-1}x)+\mathbf1\bigr)
       \in[0,1]^m,\qquad x\in Q.
 \tag{AL.8}
\]
For fixed $w$, this is continuous jointly in $(t,z,x)$. The target is a
fixed unit cube, regardless of the size of the coefficient tuple $w$.

Choose once a finite rational triangulation of $C$ compatible with $Q$
and then barycentrically subdivide. Thus $Q$ is a full subcomplex.
Write $\lambda_v(x)$ for its piecewise-affine barycentric functions and
put
\[
 \lambda(x)=\sum_{v\in Q^{(0)}}\lambda_v(x).
\]
When $\lambda(x)>0$, the point
\[
 p_Q(x)=\lambda(x)^{-1}
          \sum_{v\in Q^{(0)}}\lambda_v(x)v
\]
belongs to $Q$: the vertices with positive weights form a face, and
fullness makes that face a simplex of $Q$. Fix $c_0\in[0,1]^m$, for
example its origin, and extend (AL.8) to the cube by
\[
 E_a(x)=
 \begin{cases}
 \lambda(x)F_a(p_Q(x))+(1-\lambda(x))c_0,&\lambda(x)>0,\\
 c_0,&\lambda(x)=0.
 \end{cases}
 \tag{AL.9}
\]
This is one bounded definable formula, equals $F_a$ on $Q$, and has
values in $[0,1]^m$. For each fixed $w$, it is jointly continuous in
$(t,z,x)$. At $\lambda=0$ this follows from
$\|E_a(x)-c_0\|_\infty\le\lambda(x)$, uniformly in all parameters.
The construction neither requires a compact parameter space nor makes
any continuity assertion across the coefficient variable $w$.

\textbf{Applying the already-proved spatial-clock lemma with all coefficients retained.}\quad
Run A154's finite coordinate-variation construction on the bounded vector
\[
 G^0_a(x)=(x,E_a(x))\in[0,1]^{v+m}.
 \tag{AL.10}
\]
Its domain coordinates are the $v$ cube variables $x$; the formal
parameters are all of $(w,t,z)$. We recall exactly why the needed
uniformity is available. At a stage, every scalar component of the
current vector is a fixed definable function. O-minimal uniform
monotonicity gives an ordinary finite number of monotone intervals on
every line in the next cube coordinate, uniformly in all other
coordinates and all formal parameters. A common finite number $N_i$
works for the finitely many components at that stage.

Their partial line variations are given by the finite maximum formula
in A154: one maximizes the sum of absolute consecutive differences
over a partition with that fixed bounded number of points. The formula
uses no unbounded index set. Each component lies in $[0,1]$, so its
variation is at most $N_i$. Thus the clock
$A_a(s)=s+V_a(s)$ has its terminal value between $1$ and an ordinary
integer depending only on $N_i$ and $v+m$. The normalized inverse clock
is defined by
\[
 A_a(h_a(u))=u A_a(1),\qquad0\le u\le1.
 \tag{AL.11}
\]
The finite maximum formula and the cancellation estimate proved in
A154 give continuity and the old-coordinate Lipschitz estimates.
The inclusion of the evolving domain map among the vector coordinates
in (AL.10) gives the same estimates for that domain map itself.
The new-coordinate estimate is bounded by $A_a(1)$.

Only finitely many coordinate stages occur. Each new $N_i$ is the
uniform monotonicity count for one already fixed finite formula, so
all the final amplification bounds are ordinary finite expressions
in ordinary integers. No coefficient-dependent mesh, unbounded
iteration, or arbitrary element of the field is used as a substitute
for an ordinary finite bound. We obtain cube homeomorphisms
$\xi_a:C\to C$ such that both
\[
 \xi_a\quad\hbox{and}\quad E_a\xi_a
 \tag{AL.12}
\]
have one ordinary integer spatial Lipschitz bound, independent of all
$a=(w,t,z)$. Every coordinate change fixes the endpoints of the changed
coordinate and leaves the others unchanged. Hence it preserves every
coordinate face, so $\xi_a(Q)=Q$.

For fixed $w$, all the input functions in (AL.9) are continuous jointly
in the remaining parameters. The continuity part of A154, applied with
those parameters, proves joint continuity of $\xi_a$ and its inverse
in $(t,z,x)$. The same fixed finite formulas are used for every $w$;
continuity across different values of $w$ is not needed. Equivalently,
if such additional continuity were desired, compact graph triviality
of the bounded extension $t\Theta$ would supply finitely many
coefficient pieces, but no such partition is used in this proof.

The finite monotonicity bounds, the clock formulas, the homeomorphism
conditions, and the final displayed Lipschitz estimates are first-order
statements over the original named coefficients and the fixed rational
PL data. Their validity is also proved by the all-real-closed-field
clock argument itself. Consequently these same formulas and ordinary
integers work in every model of the fixed complete theory. No extra
unnamed coefficients have been selected. In the semialgebraic case
all of the formulas remain semialgebraic.

\textbf{The exact angular gauge and its bounds.}\quad
Set
\[
 r_{w,t,z}=b^{-1}\xi_{w,t,z}b.
 \tag{AL.13}
\]
These are homeomorphisms of $\Sigma$ with the required joint evaluation
and inverse continuity for fixed $w$. Their spatial Lipschitz bound
is the product of the standard bounds of $b,b^{-1}$ and $\xi$.
On $Q$, the extension (AL.9) equals (AL.8). Since $\xi$ preserves $Q$,
\[
 \widehat\Theta_w(t,z,u)
       =2E_{w,t,z}(\xi_{w,t,z}(b(u)))-\mathbf1.
 \tag{AL.14}
\]
This has a standard Lipschitz bound in $u$ by (AL.12). Enlarge an
ordinary integer $K$ to cover the two bounds. Finally take the maximum
over the finitely many original radial models. This proves (AL.3).
No bound on $r^{-1}$, nor any bound in the variables $z,t,w$, is asserted.

Formula (AL.5) has inverse
$\chi_{w,t}^{-1}(z,u)=(z,r_{w,t,z}^{-1}(u))$, so it is a homeomorphism
with the required joint continuity. In particular
\[
 \operatorname{pr}_L\chi_{w,t}
 =\operatorname{pr}_L\chi_{w,t}^{-1}=\operatorname{pr}_L.
\]
The substitutions (AL.6) therefore give
\[
 \widehat\Theta_w(t,\widehat A_t(z,u))
   =\Theta_w(t,A_t(z,u)),
\]
and give equality in (AL.7) at every positive time, before taking a
limit. All homeomorphism and continuity conditions in A5.019 are
preserved in both directions.

\textbf{Finite pools, not just individual existence.}\quad
Let $\gamma_j$ be one finite graph formula for an original angular
straightener, with its arbitrary finite coefficient tuple. Composition
with the one fixed formula for $\chi^{-1}$ gives one finite graph formula
$\widehat\gamma_j$. Conversely composition with $\chi$ sends any one
normalized certificate formula back to an original one. These operations
introduce only finitely many quantified variables, reuse the same
original scalar-certificate tuple $w$, and introduce no new parameter
index of unbounded length. They preserve validity and the error in
(AL.7) exactly. Thus they preserve the sentinel/liminf predicates of
A5.019 in both directions. A finite list of formulas becomes a finite
list, in every model of the same theory. Existential quantification
over scalar certificates above a centre is unaffected. This proves the
claimed equivalence of the two all-model finite-pool assertions.
\end{proof}

\paragraph{What this restriction changes and what it leaves open.}
The output class has a single ordinary integer bound in the angular
variable for the entire fixed family, including every coefficient
value and every positive radius, and the formula implementing that
normalization is fixed. The original A5.019 hypotheses impose no such
bound. For example, on the top side of the square circle one can map
$(-1+2s,1)$ to $(-1+2\sqrt{s},1)$, $0\le s\le1$, and take the identity
on the other three sides. This is a semialgebraic circle homeomorphism
with no Lipschitz constant. Its radial extension is a semialgebraic
homeomorphism of the plane and a proper product example with $L$ a
point. Thus the regularity restriction is a proper restriction on the
input angular maps; it is not a property they automatically satisfy.

The proof supplies no uniform bound for $A_t$ or $A_t^{-1}$ and no
bound on the number or complexity of their possible formula formats.
It also leaves all variation in the $L$ coordinate and in the radius
untouched. Consequently it does not establish the finite pool, an exact
angular extension theorem, or original A5. Its precise additional
reduction is to a fixed-formula family with a standard uniform angular
Lipschitz bound while retaining the standard projection control exactly.
